
Foundation model pretraining commonly uses static domain mixing ratios, such as 60\% web text, 20\% code, and 20\% books, maintained throughout training. These ratios are determined through hyperparameter search. Non-convex optimization exhibits path dependence: early gradient updates influence the representational geometry available for subsequent learning. This raises the question of whether temporal ordering of domains matters.

Standard practice treats temporal composition implicitly. Methods like DoReMi optimize static mixing ratios through group distributionally robust optimization, while two-phase training (general pretraining followed by domain-specific continuation) implements a single sharp transition. These approaches determine how much data from each domain to include but largely ignore when to present it.

This paper reports proof-of-concept implementation of diversity-ranked domain scheduling. The approach orders training domains from high to low diversity during pretraining using smooth Gaussian-weighted transitions. Diversity is quantified through corpus statistics: vocabulary entropy, syntactic complexity, and semantic spread. The hypothesis is that early high-diversity exposure establishes broader gradient covariance geometry through path-dependent optimization.

The implementation uses GPT-2 style transformer models (1B and 7B scale) with The Pile dataset. Six domains (Pile-CC, StackExchange, Wikipedia, ArXiv, Github, PubMed Central) are ranked by diversity scores (0.92 to 0.35). Four experimental conditions are implemented: static mixture (uniform sampling), diversity-ranked (high-to-low), reversed (low-to-high), and shuffled (randomized order).

\textbf{This is a proof-of-concept validation only.} All core components execute without errors. The curriculum scheduler produces valid domain sampling probabilities. Unit tests (22 total) pass. A smoke test (10 training steps, 1B scale) completes successfully. However, the smoke test does not demonstrate convergence, statistical significance, or performance improvement. The composite benchmark score of 0.2558 reflects an untrained model trained for only 10 steps.

Full-scale validation is deferred. Testing the performance hypothesis (≥2.0\% improvement at 1B, ≥0.5\% at 7B) requires training to convergence (100,000-150,000 steps) with multiple seeds (n=5) for statistical testing. The mechanistic explanation (gradient covariance → representational persistence → robust learning) requires measurements (Participation Ratio at multiple checkpoints, CKA similarity between checkpoint pairs, Fisher overlap during continual learning) not performed in this work.

The remainder of this paper is organized as follows. Section 2 reviews related work. Section 3 describes the diversity-ranked scheduling methodology and proof-of-concept protocol. Section 4 details experimental setup. Section 5 reports proof-of-concept results. Section 6 discusses limitations and planned full-scale experiments. Section 7 concludes.
